% TOP_PROBLEM: {"id": 3093, "title": "Convex Equipartitions with Extreme Perimeter", "category_id": 6, "category": {"id": 6, "name": "geometry", "display_name": "Geometry", "description": "Euclidean and non-Euclidean geometry, geometric structures.", "slug": "geometry", "order_index": 6, "created_at": "2026-07-31T15:26:25.671Z"}, "status": "open", "problem_number": "OPG-60010", "set_id": 12, "statusQualification": "The region is specified as a compact convex body with interior, common boundaries are counted in each adjacent perimeter, and extrema are defined through infimum and supremum before attainment."}
% FIRST_POSTED: 2026-10-01
% ATTEMPT_STATUS: unresolved
% COMPLETION: 5
% MODEL: GPT 6 Astra Ultra
% UnsolvedMath ID 3093; source catalogue code OPG-60010
% !TeX program = lualatex
% Research attempt, 2026-10-01. Canonical statement preserved verbatim.
\documentclass[11pt,a4paper]{article}
\usepackage[margin=25mm]{geometry}
\usepackage{fontspec}
\setmainfont{lmroman10-regular.otf}[BoldFont=lmroman10-bold.otf,
 ItalicFont=lmroman10-italic.otf,BoldItalicFont=lmroman10-bolditalic.otf]
\usepackage{amsmath,amssymb,mathtools}
\usepackage{unicode-math}
\setmathfont{latinmodern-math.otf}
\AtBeginDocument{\let\setminus\smallsetminus}
\usepackage{xurl,hyperref}
\hypersetup{colorlinks=true,urlcolor=blue}
\setcounter{secnumdepth}{0}
\setlength{\parindent}{0pt}
\setlength{\parskip}{5pt}
\setlength{\emergencystretch}{3em}
\begin{document}
\section*{Convex Equipartitions with Extreme Perimeter}\label{3093}
{\small UnsolvedMath ID 3093 · OPG-60010 · Geometry · OpenGarden}
\subsection{Definitions and mathematical statement}
Let $C\subset\mathbb R^2$ be a compact convex set with nonempty interior,
and let $n\ge1$ be an integer. Write $|C|$ for its area and
$\operatorname{per}(C)$ for the Euclidean length of its boundary.
An admissible convex equipartition is a family of compact convex sets
$C_1,\ldots,C_n$ whose union is $C$, whose interiors are pairwise disjoint,
and for which $|C_i|=|C|/n$ for every $i$.

The quantity to optimize is the sum
\[
 P(C_1,\ldots,C_n)=\sum_{i=1}^n\operatorname{per}(C_i).
\]
Determine its least and greatest possible values over admissible partitions,
and describe partitions attaining them. Equivalently one can first define
$P_{\min}(C,n)$ and $P_{\max}(C,n)$ as the infimum and supremum and
include attainment in a complete determination. For a polygon the pieces
are understood as closed regions meeting along their cut boundaries,
with boundary overlaps allowed but no overlap of positive area.

The exterior boundary of $C$ is counted once in the sum, while every
internal interface of positive length is counted twice. Thus, writing
$L$ for the total length of internal interfaces counted once,
$P=\operatorname{per}(C)+2L$. The optimization is also optimization of
cut length. Only areas must be equal; piece perimeters may differ and
forcing them to coincide would change the feasible family.

The question ranges over the input body and number of pieces, including
nonsymmetric bodies. Merely establishing that equal-area convex partitions
exist does not identify either perimeter extremum.
\subsection{Short English statement}
Among convex partitions into equally large areas, which arrangements minimize or maximize the sum of piece perimeters?
\subsection{Research attempt}
\paragraph{Scope and literature check.}
The source question [S2], posted by R. Nandakumar and checked 1 October 2026,
asks about arbitrary convex bodies and all numbers of pieces. Its proposed
greedy maximization is a conjectural route, not a theorem used here. This
attempt completely computes the two-piece extrema for centrally symmetric
bodies, with explicit formulas for rectangles, and identifies why that
calculation supplies no general greedy rule.

\paragraph{Two convex pieces force one straight cut.}
Suppose $C=C_1\cup C_2$ is an admissible two-piece partition. Both pieces have
positive area. Their disjoint convex interiors can be separated by a line;
this is the elementary hyperplane separation theorem in dimension two [A1].
Take the separating closed half-planes $H_1,H_2$ so that $C_i\subset H_i$.
Every point of $C$ strictly inside $H_i$ must then belong to $C_i$, since it
cannot belong to the other piece. Closedness gives $C_i=C\cap H_i$, including
the separating chord. Thus every admissible partition is a line bisection,
and if that chord has length $l$ the perimeter sum is
\[
                         P=\operatorname{per}(C)+2l.
\]
The two copies of the chord, not one, must be counted.

\paragraph{Central symmetry fixes the cut's position.}
Translate the centre of symmetry to the origin. For a fixed unit normal $u$
let $A_u(t)$ be the area of $C\cap\{x:\langle x,u\rangle\le t\}$.
It is continuous in $t$, since the area of a strip tends to zero with its
width. It is strictly increasing between the supporting lines: any nonempty
open strip there meets the interior of $C$ in a set of positive area.
Symmetry gives $A_u(0)=|C|/2$. Therefore zero is the unique bisecting offset,
and every equal-area two-piece cut passes through the origin.
Let $r(v)>0$ be the distance from the origin to the boundary along a unit
direction $v$. The corresponding central chord has length $2r(v)$.
Consequently
\[
 P_{\min}(C,2)=\operatorname{per}(C)+4\min_{v\in S^1}r(v),\qquad
 P_{\max}(C,2)=\operatorname{per}(C)+4\max_{v\in S^1}r(v).
\]
Both extrema are attained. Indeed the radial function is continuous for a
compact convex body containing the origin in its interior, and $S^1$ is
compact. Continuity also follows directly from the continuous Minkowski
functional of such a body [A1].

\paragraph{Explicit rectangle and disk certificates.}
For $C=[-a,a]\times[-b,b]$ with $a,b>0$, a central chord in direction
$(\cos\theta,\sin\theta)$ has half-length
\[
 r(\theta)=\min\{a/|\cos\theta|,b/|\sin\theta|\},
\]
with division by zero interpreted as infinity. The smallest half-length is
$\min(a,b)$: the disk of that radius is contained in the rectangle, and a
normal to a shortest pair of opposite sides attains it. The largest is
$\sqrt{a^2+b^2}$, attained in the corner direction and bounded above by the
corner distance for every point of the rectangle. Hence
\[
 P_{\min}=4(a+b)+4\min(a,b),\qquad
 P_{\max}=4(a+b)+4\sqrt{a^2+b^2}.
\]
For a disk of radius $R$, all half-lengths equal $R$ and every diameter cut
has perimeter sum $2\pi R+4R$.

\paragraph{Verification and obstruction to iteration.}
The area of each central half is exactly half by the symmetry map $x\mapsto-x$;
no assumption that their perimeters agree was needed. For a square the
minimum cut is parallel to a side and the maximum is a diagonal, consistent
with their actual chord lengths. When $n>2$, the first piece has area
$|C|/n$, so its cut generally does not pass through the centre; moreover the
remainder usually loses central symmetry. Maximizing the first chord need
not maximize the sum of later chord lengths. Neither the reduction to one
line nor the displayed radial formula justifies the source's greedy proposal.

\paragraph{Final status.}
\textbf{UNRESOLVED.} Both extrema are determined for two pieces of every
centrally symmetric planar convex body. The requested arbitrary-body,
arbitrary-$n$ optimization is not determined. Completion estimate: 5\%.

\subsection{Sources}
\begin{itemize}
\item[S1] ULAM.AI and the UnsolvedMath Contributors, supplied problems.json, record 3093 (OPG-60010), OpenGarden collection. Catalogue identity and source transcription; CC BY 4.0.
\url{https://huggingface.co/datasets/ulamai/UnsolvedMath}.
\item[S2] Open Problem Garden, Convex Equipartitions with Extreme Perimeter. Original problem and discussion; the mathematical formulation below is expanded from the supplied source record.
\url{http://www.openproblemgarden.org/op/convex_equipartitions_with_extreme_perimeter}.

\item[A1] R. Schneider, \emph{Convex Bodies: The Brunn--Minkowski Theory}, second expanded edition, Chapter 1, Basic convexity (separation and convex-body functions). Publisher chapter record checked 1 October 2026. \url{https://www.cambridge.org/core/books/abs/convex-bodies-the-brunnminkowski-theory/basic-convexity/A9304A212E8B47356C5514FD9CC217CD}.
\end{itemize}
\end{document}
